\subsection{Theoretical Background}
\label{sec:theoretical_background}

\subsubsection{Traffic Video Detection}
\label{sec:traffic_video_detection}

Intelligent transport systems use visual sensing to estimate vehicle flows, classify vehicle types, and support traffic management decisions. Camera-based monitoring is attractive because existing surveillance infrastructure can capture dense spatial information without installing physical sensors at every lane. At the same time, urban video presents perceptual variability that is uncommon in controlled image benchmarks: shadows, night and daylight transitions, camera height differences, occlusion, congestion, and heterogeneous vehicle shapes can all change the appearance of the same semantic class~\cite{Zhou2025TrafficSurveillanceSurvey,Azfar2024ComplexTrafficReview}. These factors motivate methods that treat dataset preparation, annotation geometry, and evaluation design as part of the detection problem rather than as secondary implementation details.

\subsubsection{Oriented Bounding Boxes}
\label{sec:oriented_bounding_boxes}

Object detectors commonly localize instances with horizontal bounding boxes (HBBs), but HBBs are inefficient for vehicles observed at arbitrary headings. A diagonal car or turning bus can occupy only a fraction of its axis-aligned rectangle, causing extra background pixels to enter the localization target and increasing overlap between neighboring vehicles. Oriented bounding boxes (OBBs) add an angle parameter to the center and size representation, or equivalently encode the four rotated box corners, so the box can align with the physical vehicle footprint~\cite{Wang2025OBBSurvey,Zand2022OBB}. OBBs are therefore well suited to elevated traffic cameras and aerial viewpoints, where vehicle orientation is informative and localization errors can propagate to safety-related downstream quantities~\cite{Ahmed2025}.

\subsubsection{Rotated Intersection over Union}
\label{sec:riou_metric}

Rotated Intersection over Union (rIoU) measures the geometric agreement between two oriented boxes after converting them into polygons. Unlike conventional IoU for horizontal rectangles, rIoU computes the intersection and union areas of the rotated quadrilaterals themselves. It is therefore sensitive to angular misalignment, not only to center displacement or scale error. This distinction is important for elongated vehicles: two boxes can have similar centers and dimensions but a poor rIoU if their headings differ substantially. In this work, the metric implementation validates every OBB, converts it to a four-vertex polygon, computes convex-polygon overlap with OpenCV, and clips numerical results to the interval $[0,1]$. For a predicted box $B_p$ and a ground-truth box $B_g$, rIoU is defined as
\begin{equation}
  \operatorname{rIoU}(B_p,B_g) =
  \frac{\left|P(B_p)\cap P(B_g)\right|}
       {\left|P(B_p)\cup P(B_g)\right|},
  \label{eq:riou}
\end{equation}
where $P(\cdot)$ denotes the polygon induced by an OBB and $|\cdot|$ denotes polygon area. A value of 1 indicates perfect overlap, while 0 indicates no overlap. Values near a decision threshold, such as 0.50 or 0.80, are interpreted as localization matches only if the rotated geometry satisfies that threshold.

\subsubsection{Macro AP-rIoU}
\label{sec:macro_ap_riou}

The evaluation protocol implements Macro AP-rIoU over the nine vehicle classes and seven rIoU thresholds, $\mathcal{T}=\{0.50,0.55,\ldots,0.80\}$. For each class and threshold, predictions are sorted by confidence and greedily matched to unused ground-truth boxes from the same frame and class. Average Precision (AP) is then computed with 101 interpolated recall levels. The final score gives equal weight to every class, so rare classes affect the metric as much as frequent ones:
\begin{equation}
\begin{aligned}
  \operatorname{MacroAP}_{rIoU}
  = \frac{1}{9}\sum_{c=1}^{9}
    \frac{1}{|\mathcal{T}|}\sum_{\tau\in\mathcal{T}}
    \operatorname{AP}_{c,\tau}.
\end{aligned}
\label{eq:macro_ap_riou}
\end{equation}
Macro averaging is important in a long-tailed dataset. A detector that performs well on cars but poorly on microbuses, articulated buses, or mototaxis can still obtain a limited final score because class frequency does not weight the average. The metric therefore supports the methodological focus on minority-class data augmentation.

\subsubsection{Small Objects and Class Imbalance}
\label{sec:small_objects_imbalance}

Traffic scenes from elevated cameras contain many compact or distant vehicles that occupy few pixels. Small-object surveys identify repeated downsampling, limited feature support, and dense context as common failure modes for convolutional detectors~\cite{Cheng2023SmallObjectSurvey,Nikouei2025SmallObjectSurvey}. Class imbalance compounds the problem because frequent classes dominate the training signal. In the study dataset, the car class accounts for most annotations, while several bus and mototaxi categories occupy the long tail. Prior work on motorcycle detection and oriented aerial vehicle detection motivates both class-aware evaluation and data-level balancing strategies~\cite{Espinosa2020,Mo2020}. The augmentation strategy follows that motivation by allocating synthetic additions only to classes whose real frequency falls below a configured threshold.

\subsubsection{Segmentation-Assisted Synthetic Augmentation}
\label{sec:segmentation_assisted_augmentation_background}

Synthetic data can improve detector training when it adds useful diversity without destroying label validity. Artificial vehicle images and broader surveys show that combining real and synthetic data may be effective, but they also highlight the risk of domain gaps and unrealistic artifacts~\cite{Weber2021ArtificialImagesVehicleDetection,Mumuni2024}. More recent generative augmentation approaches use segmentation or region control to preserve object identity and spatial annotations, because unconditioned generation can shift object boundaries or create label-image mismatch~\cite{Rahat2024,Benkedadra2024,ZhuH2025}. The method documented in this paper takes a conservative form of this idea: it uses segmentation masks for foreground extraction and copy-paste composition, while avoiding full-frame diffusion generation in the final training set.
